% 08-build.tex
\chapter{Development and Build System}
\label{ch:build}

\section{Build Process}
\label{sec:build-process}

Doki uses a Make-based build system that compiles four binaries for
multiple architectures. The build produces statically linked executables
with no runtime CGo dependencies.

\subsection{Compilation Targets}
\label{sec:compilation-targets}

\begin{figure}[htbp]
  \centering
  \begin{tikzpicture}[
    node distance=0.5cm,
    every node/.style={font=\sffamily\footnotesize},
    stage/.style={
      rectangle, rounded corners=2pt,
      draw=nodeblue!60, fill=fillblue,
      minimum width=2.2cm, minimum height=0.7cm,
      text=textdark, align=center, inner sep=4pt
    }
  ]
    \node[stage] (src) {Go Source\\{\scriptsize\color{textgray} cmd/\& pkg/}};
    \node[stage, right=0.8cm of src] (build) {Make\\{\scriptsize\color{textgray} build target}};
    \node[stage, right=0.8cm of build] (cross) {Cross-\\{\scriptsize\color{textgray} compile}};
    \node[stage, right=0.8cm of cross] (strip) {Strip\\{\scriptsize\color{textgray} -ldflags "-s -w"}};
    \node[stage, right=0.8cm of strip] (archive) {Archive\\{\scriptsize\color{textgray} .tar.gz}};

    \draw[arr] (src) -- (build);
    \draw[arr] (build) -- (cross);
    \draw[arr] (cross) -- (strip);
    \draw[arr] (strip) -- (archive);

    % Cache hit
    \draw[arrDash, nodeorange] (build.south) -- ++(0,-0.5)
      node[below, font=\sffamily\tiny, text=nodeorange] {incremental}
      -- ++(2,0) |- (strip.south);
  \end{tikzpicture}
  \caption{Build pipeline from source to release archive.}
  \label{fig:build-pipeline}
\end{figure}

\begin{table}[htbp]
  \centering
  \caption{Build targets and their output binaries.}
  \label{tab:build-targets}
  \begin{tabularx}{\textwidth}{l l X}
    \toprule
    \textbf{Target} & \textbf{Output} & \textbf{Description} \\
    \midrule
    \texttt{build}            & \texttt{bin/doki*}         & Development binaries \\
    \texttt{build-android-arm64}  & \texttt{releases/doki-v*-arm64.tar.gz}  & Android ARM64 \\
    \texttt{build-android-armv7}  & \texttt{releases/doki-v*-armv7.tar.gz}  & Android ARMv7 \\
    \texttt{build-linux-arm64}    & \texttt{releases/doki-v*-arm64.tar.gz}  & Linux ARM64 \\
    \texttt{build-linux-armv7}    & \texttt{releases/doki-v*-armv7.tar.gz}  & Linux ARMv7 \\
    \texttt{build-darwin-arm64}   & \texttt{releases/doki-v*-arm64.tar.gz}  & macOS ARM64 \\
    \bottomrule
  \end{tabularx}
\end{table}

\subsection{Android ARMv7 Build Fix}
\label{sec:armv7-fix}

A critical build issue was resolved in v0.9.3: Go 1.22+ requires
an external linker for \texttt{GOOS=android GOARCH=arm}, causing
a ``requires external cgo linking'' error with \texttt{CGO\_ENABLED=0}.

The solution builds ARMv7 with \texttt{GOOS=linux} instead of
\texttt{GOOS=android}. Doki is 100\% CGo-free, and Android detection
uses filesystem probes (\texttt{/data/data/com.termux}) rather than
\texttt{runtime.GOOS}. The resulting binaries compile as ELF 32-bit
ARM EABI5 statically linked executables.

\begin{codebox}[title={Makefile fix for ARMv7}]{makefile}
build-android-armv7:
	CGO\_ENABLED=0 GOOS=linux GOARCH=arm GOARM=7 \\
	  go build -ldflags "-s -w" \\
	  -o bin/doki ./cmd/doki/
\end{codebox}

\section{Dokifile Build System}
\label{sec:dokifile}

The Dokifile build system parses Dockerfile instructions and produces
OCI-compliant images. It supports multi-stage builds, build arguments,
and a layer cache for incremental rebuilds.

\subsection{Build Pipeline}
\label{sec:build-pipeline}

\begin{figure}[htbp]
  \centering
  \begin{tikzpicture}[
    node distance=0.5cm,
    every node/.style={font=\sffamily\footnotesize},
    stage/.style={
      rectangle, rounded corners=2pt,
      draw=nodeblue!60, fill=fillblue,
      minimum width=2.2cm, minimum height=0.8cm,
      text=textdark, align=center, inner sep=4pt
    }
  ]
    \node[stage] (parse) {Parse\\Dokifile};
    \node[stage, right=0.8cm of parse] (resolve) {Resolve\\Base Image};
    \node[stage, right=0.8cm of resolve] (cache) {Check\\Layer Cache};
    \node[stage, right=0.8cm of cache] (exec) {Execute\\Instructions};
    \node[stage, right=0.8cm of exec] (commit) {Commit\\Layers};
    \node[stage, right=0.8cm of commit] (config) {Apply\\Config};

    \draw[arr] (parse) -- (resolve);
    \draw[arr] (resolve) -- (cache);
    \draw[arr] (cache) -- (exec);
    \draw[arr] (exec) -- (commit);
    \draw[arr] (commit) -- (config);

    % Cache hit arrow
    \draw[arrDash, nodeorange] (cache.south) -- ++(0,-0.6)
      node[below, font=\sffamily\tiny, text=nodeorange]
        {cache hit: skip to commit}
      -- ++(1.5,0) |- (commit.south);
  \end{tikzpicture}
  \caption{Dokifile build pipeline with layer cache optimization.}
  \label{fig:build-pipeline}
\end{figure}

\section{Compose Orchestration}
\label{sec:compose-orchestration}

The Compose engine processes YAML configuration files and manages
multi-container applications. It supports service dependencies,
health checks, profiles, and volume definitions.

\subsection{Compose File Structure}
\label{sec:compose-structure}

A typical Compose file defines services, networks, and volumes:

\begin{codebox}[title={Example Compose file}]{yaml}
services:
  web:
    build: ./web
    ports:
      - "8080:80"
    depends\_on:
      api:
        condition: service\_started
  api:
    image: myapi:latest
    environment:
      - DB\_HOST=db
    depends\_on:
      db:
        condition: service\_healthy
  db:
    image: postgres:16
    healthcheck:
      test: ["CMD", "pg\_isready"]
      interval: 10s
volumes:
  pgdata:
\end{codebox}

\subsection{Service Dependency Resolution}
\label{sec:dependency-resolution}

The engine resolves dependencies using topological sort. Services
without dependencies start first, and dependent services wait for
their prerequisites to reach the specified condition.

\begin{figure}[htbp]
  \centering
  \begin{tikzpicture}[
    node distance=0.8cm,
    every node/.style={font=\sffamily\footnotesize},
    svc/.style={
      rectangle, rounded corners=2pt,
      draw=nodeblue!60, fill=fillblue,
      minimum width=1.8cm, minimum height=0.7cm,
      text=textdark, align=center, inner sep=4pt
    },
    arr/.style={
      -{Stealth[length=2.5mm, width=2mm]},
      linegray, line width=0.5pt
    }
  ]
    % Services
    \node[svc] (db) {db};
    \node[svc, right=1.5cm of db] (redis) {redis};
    \node[svc, above=1.0cm of db] (api) {api};
    \node[svc, right=1.5cm of api] (worker) {worker};
    \node[svc, above=1.0cm of api] (web) {web};
    \node[svc, right=1.5cm of web] (nginx) {nginx};

    % Dependencies
    \draw[arr] (api) -- (db);
    \draw[arr] (api) -- (redis);
    \draw[arr] (worker) -- (db);
    \draw[arr] (worker) -- (redis);
    \draw[arr] (web) -- (api);
    \draw[arr] (nginx) -- (web);

    % Start order
    \node[annot, below=0.3cm of db] {1. db, redis};
    \node[annot, below=0.3cm of api] {2. api, worker};
    \node[annot, below=0.3cm of web] {3. web, nginx};
  \end{tikzpicture}
  \caption{Compose service dependency graph with start order.}
  \label{fig:compose-dependencies}
\end{figure}
